% G3-2 Pengumpulan data.
% Sumber: diolah penulis; seluruh rincian dari
% context/konteks_collector_data_v3.0.md (Bagian 1, 3, 4, 6).
\begin{tikzpicture}[x=1mm,y=-1mm,font=\footnotesize,
    sumber/.style={draw,rectangle,rounded corners=2pt,text width=42mm,
      align=center,inner sep=2.5pt,minimum height=13mm},
    proses/.style={draw,rectangle,text width=28mm,align=center,
      inner sep=2.5pt,minimum height=13mm},
    tabel/.style={draw,rectangle,text width=34mm,align=center,
      inner sep=2.5pt,minimum height=15mm},
    takpakai/.style={tabel,dashed},
    berkas/.style={draw,trapezium,trapezium left angle=72,
      trapezium right angle=108,trapezium stretches body,
      text width=28mm,align=center,inner sep=2.5pt,minimum height=13mm},
    panah/.style={-{Stealth[length=2.2mm]}}]

  % ===== kolom 1: API Indodax, dua endpoint =====
  \node[anchor=south,font=\small,align=center] at (24,-10)
    {\textbf{API publik Indodax}};
  \node[sumber] (ep1) at (24,4)
    {\textit{endpoint} ticker\\\texttt{/api/btc\_idr/ticker}};
  \node[sumber] (ep2) at (24,32)
    {\textit{endpoint candle}\\\texttt{/tradingview/}\\\texttt{history\_v2}, \texttt{tf=1}};

  % ===== kolom 2: collector =====
  \node[proses] (col) at (66,18)
    {\textit{Collector} (Python)\\di VPS\\\textit{cron} tiap 1 menit, UTC};
  \draw[panah] (ep1.east) -- (48,4)  -- (48,18) -- (col.west);
  \draw[panah] (ep2.east) -- (48,32) -- (48,18) -- (col.west);

  % ===== kolom 3: dua tabel SQLite =====
  \node[anchor=south,font=\small,align=center] at (108,-10)
    {\textbf{Basis data SQLite}};
  \node[takpakai] (tb1) at (108,2)
    {\texttt{ticker\_snapshots}\\\textit{last}, \textit{bid}, \textit{ask},
     data 24 jam};
  \node[anchor=north,align=center,text width=40mm,font=\footnotesize]
    at (108,12.5) {disimpan, tidak dipakai\\dalam eksperimen};
  \node[tabel] (tb2) at (108,36)
    {\texttt{ohlc\_1m}\\\textit{open}, \textit{high}, \textit{low},
     \textit{close}, volume};
  \draw[panah,dashed] (col.east) -- (88,18) -- (88,2)  -- (tb1.west);
  \draw[panah]        (col.east) -- (88,18) -- (88,36) -- (tb2.west);

  % ===== baris bawah: dua dataset eksperimen =====
  \node[berkas] (ds1) at (20,74) {\textbf{dataset}\\\textit{candle} 1 menit};
  \node[proses,text width=38mm] (agg) at (66,74)
    {\textbf{agregasi 5 menit}\\\textit{open} pertama, \textit{high} maks,
     \textit{low} min, \textit{close} terakhir, volume dijumlah};
  \node[berkas] (ds2) at (112,74) {\textbf{dataset}\\\textit{candle} 5 menit};

  % dari ohlc_1m turun ke baris dataset (garis di y=54, di atas kotak agregasi)
  \draw (108,44) -- (108,56) -- (20,56);
  \draw[panah] (20,56) -- (ds1.north);
  \draw[panah] (ds1.east) -- (agg.west);
  \draw[panah] (agg.east) -- (ds2.west);

\end{tikzpicture}
